Natural Language Inference (NLI) \cite{fyodorov2000natural} concerns the ability to understand the relationship between two sentences. In practice, this task is usually structured as a two or three class classification problem where the model classifies whether the second sentence logically follows from the first, contradicts the first sentence, or is possibly true (neutral). SuperGLUE includes an NLI dataset, RTE, which evaluates the binary version of the task. On RTE, only the largest version of GPT-3 performs convincingly better than random (56\%) in any evaluation setting, but in a few-shot setting GPT-3 performs similarly to a single-task fine-tuned BERT Large. We also evaluate on the recently introduced Adversarial Natural Language Inference (ANLI) dataset \cite{nie2019adversarial}. ANLI is a difficult dataset employing a series of adversarially mined natural language inference questions in three rounds (R1, R2, and R3). Similar to RTE, all of our models smaller than GPT-3 perform at almost exactly random chance on ANLI, even in the few-shot setting ($\sim33\%$), whereas GPT-3 itself shows signs of life on Round 3. Results for ANLI R3 are highlighted in Figure \ref{graph:anli} and full results for all rounds can be found in Appendix \ref{appendix:results_on_all_tasks}. These results on both RTE and ANLI suggest that NLI is still a very difficult task for language models and they are only just beginning to show signs of progress.